\documentclass[10pt,a4paper]{amsart}
\usepackage[margin=19mm]{geometry}
\usepackage{amsmath,amssymb,mathtools}
\usepackage[scr=rsfs,cal=euler]{mathalfa}
\usepackage{xfrac}
\usepackage[unicode,hidelinks,bookmarks=false]{hyperref}
\newcommand{\ie}{i.e.\ }
\newcommand{\cf}{cf.\ }
\newcommand{\resp}{respectively\ }
\newcommand{\Subsection}{\subsection}
\providecommand{\nobreakdash}{\nobreak\hbox{-}}
\setlength{\emergencystretch}{2em}
\multlinegap0pt
\usepackage{bm}
\usepackage{bbm}
\usepackage{dsfont}
\usepackage{xspace}
\usepackage{cancel}
\usepackage{mathtools}
\usepackage{amsthm}
\newtheorem{lemma}{Lemma}
\newtheorem{proposition}{Proposition}
\title[Complete characterization of $2$-near perfect numbers with e]{Complete characterization of $2$-near perfect numbers with exactly 2 prime factors}
\author{Richard Fearon, Henry Foushee, Benjamin Porosoff, Alexander Skula, Joshua Zelinsky, Kyle Zhang}
\date{Source: arXiv:2508.06651v3 (CC BY 4.0)}
\begin{document}
\maketitle

\noindent\emph{Source and licence.} Complete characterization of $2$-near perfect numbers with exactly 2 prime factors. Version arXiv:2508.06651v3 (CC BY 4.0), distributed under the Creative Commons Attribution 4.0 International licence, as recorded in the arXiv licence metadata for this exact version and shown on the abstract page https://arxiv.org/abs/2508.06651v3. Full rights notice: licenses/arxiv-2508.06651-CC-BY-4.0.txt.\par\smallskip

\noindent\textit{Editorial scope.} Counted unit: the Proposition whose source label is \texttt{prop:no\_pm\_divisor} in the article (source0.tex lines 816--818, source label \texttt{prop:no\_pm\_divisor}), delivered together with the article's own complete original proof of it (source0.tex lines 821--861, one \texttt{proof} environment of 41 lines). No mathematics is written, repaired or re-derived: statement and proof are the article's own text line for line. Required context retained uncounted: lem:basic\_equation (lines 312--315); lem:equivalent\_form (lines 322--325); lem:parity\_constraint (lines 340--342); prop:p\_divisibility\_constraint (lines 774--776). Every definition, display and label of those lines is retained unchanged. Omitted from this package and not redistributed: 1--311, 316--321, 326--339, 343--773, 777--815, 819--820, 862--1303. Nothing inside a retained excerpt is omitted or altered: every retained body is delivered exactly as the article's lines are written, so no omission receipt of any kind accompanies this package. Citations: the retained excerpts carry no citation command at all, so no bibliography entry of the article is transcribed and no reference is redistributed. Reference adaptation: none was needed and none was made; every reference inside the delivered unit resolves inside it, so no unresolved reference or citation is produced. Printed slips kept literal: no printed word, symbol, hypothesis or number of the counted statement or of its proof was altered. Layout and class: the article's own class is \texttt{article} with packages including amsthm, amsfonts, amssymb, amsmath, epsf, verbatim, url, placeins, geometry, natbib, caption, graphicx; the wrapper is the lane's pinned \texttt{amsart} editorial wrapper at 10pt on A4, which restates the theorem-like environments the retained text needs and adds the article's own macro definitions that the retained text needs (none), with extra wrapper packages (bm, bbm, dsfont, xspace, cancel, mathtools). The delivered numbering is the unit's own. Assets: no retained span contains a figure environment, a graphics-inclusion command, a PDF-page inclusion, a table or a TikZ picture; no image file is copied into this package, the package declares no asset dependency and no image file is redistributed with it. Licence: version arXiv:2508.06651v3 of this article is distributed under the Creative Commons Attribution 4.0 International licence, as recorded in the arXiv licence metadata for this exact version and shown on the abstract page https://arxiv.org/abs/2508.06651v3. A full rights notice is delivered at licenses/arxiv-2508.06651-CC-BY-4.0.txt. Characters: every retained excerpt is pure ASCII, so the delivered engine, XeLaTeX with its default Latin Modern text font, reports no missing character.
\begin{lemma}\label{lem:basic_equation}
    Let $n=2^kp^m$ be a $2$-near perfect number with omitted divisors $d_1$ and $d_2$. Then
    \[(2^{k+1}-1)(1+p+\ldots+p^m)=2^{k+1}p^m+d_1+d_2.\]
\end{lemma}

\begin{lemma}\label{lem:equivalent_form}
    Let $n=2^kp^m$ be a $2$-near perfect number with omitted divisors $d_1$ and $d_2$. Then 
    \[(2^{k+1}-p)(1+p+\ldots+p^{m-1})=1+d_1+d_2.\]
\end{lemma}

\begin{lemma}\label{lem:parity_constraint}
    Let $n=2^kp^m$ be a $2$-near number perfect with omitted divisors $d_1$ and $d_2$. If $m$ is odd, then $d_1$ and $d_2$ have the same parity. If $m$ is even, then $d_1$ and $d_2$ have opposite parities.
\end{lemma}

\begin{proposition}\label{prop:p_divisibility_constraint}
   If $n=2^kp^m$ is $2$-near perfect, where $m\ge3$ and $2^{k+1}<p^2+1$, then $p^{m-2}$ divides one omitted divisor, and $p^2$ does not divide the other. 
\end{proposition}

\begin{proposition}\label{prop:no_pm_divisor}
    If $n=2^kp^m$ is $2$-near perfect, and $m\ge3$ and $2^{k+1}<p^2+1$, then $2^cp^m$ is not an omitted divisor.
\end{proposition}

\begin{proof}
Suppose the omitted divisors are $2^ap^b$ and $2^cp^m$.
Proposition~\ref{prop:p_divisibility_constraint} gives $b\in\{0,1\}$.
Let $S=1+p+\cdots+p^{m-1}$. Lemma~\ref{lem:equivalent_form}
and $p^m\equiv1\pmod S$ give
\[
S\mid1+2^ap^b+2^c,
\qquad S\le1+2^ap^b+2^c.
\]
If $b=0$, then $p+p^2\le S-1\le2^a+2^c\le2^{k+1}<p^2+1$,
which is impossible. Thus $b=1$.
If $m\ge4$, then
\[
p+p^2+p^3\le S-1\le2^k(p+1)
<\frac{p^3+p^2+p+1}{2}<p+p^2+p^3,
\]
also impossible. Therefore $m=3$. Lemma~\ref{lem:basic_equation} now reads
\[
(2^{k+1}-1)(p+1)(p^2+1)
=p\bigl((2^{k+1}+2^c)p^2+2^a\bigr).
\]
Since $\gcd(p,p^2+1)=1$, reduction modulo $p^2+1$ gives
\[
p^2+1\mid D,\qquad D=2^{k+1}+2^c-2^a>0.
\]
If $c\le a$, then $D\le2^{k+1}<p^2+1$, a contradiction.
Hence $c>a$. Since $m$ is odd, Lemma~\ref{lem:parity_constraint}
forces both omitted divisors to be even; thus $a\ge1$ and $c\ge2$.
Moreover,
\[
D<2^{k+1}+2^k<\frac32(p^2+1).
\]
Being a positive multiple of $p^2+1$, $D$ must equal $p^2+1$.
Modulo $4$, this equality gives $2\equiv-2^a$, hence $a=1$.
But then the left side of Lemma~\ref{lem:basic_equation} is divisible
by $4$, since $(p+1)(p^2+1)$ is, whereas its right side is
\[
2^{k+1}p^3+2p+2^cp^3\equiv2\pmod4.
\]
This contradiction excludes the proposed divisor.
\end{proof}

\end{document}
